\section{Related Work}
\textbf{Sub-Gaussian contextual bandits.} LinUCB and its variants \citep{lindqvist2019linucb} construct confidence ellipsoids around a ridge-regression estimate and select actions optimistically; their regret guarantees assume sub-Gaussian noise. Thompson sampling \citep{chen2020thompson} instead maintains a Bayesian posterior over the reward parameters and samples an action according to its posterior probability of optimality; empirically it often outperforms UCB-style methods, but existing analyses again lean on Gaussian or sub-Gaussian likelihoods.

\textbf{Heavy-tailed bandits.} For the non-contextual multi-armed setting, robust mean estimators --- truncation, median-of-means, and Catoni's estimator --- have been used to build UCB-style algorithms with regret guarantees under only finite $(1+\epsilon)$ moments \citep{holt2021heavytail, rostova2021ucb, gomez2018exploration}. \citet{tanaka2023robust} extend median-of-means to the linear contextual case, obtaining an optimistic algorithm with robust confidence sets. \citet{anand2024moments} give sharper finite-sample concentration bounds for truncated estimators that we use directly in our analysis. To our knowledge, no prior work adapts randomized (Thompson-sampling-style) exploration, rather than optimism, to the heavy-tailed contextual setting.

\textbf{Calibration in Bayesian bandits.} \citet{osei2022calibration} study posterior calibration for contextual bandits under model misspecification, showing that miscalibrated posterior variance can inflate regret even when the mean estimate is consistent. Our online moment-calibration step is inspired by this observation, but targets the distinct failure mode of heavy tails rather than model misspecification. \citet{sterling2020epsilon} and \citet{davis2017regret} provide complementary lower-bound and minimax analyses that we build on in Section~\ref{sec:method}.
